% formato-detallado.tex
% Esqueleto para escribir tesis en la FCFM.
% Actualización del original fcfmskel.tex
% --- ATENCIÓN--
% Hay varias instrucciones comentadas. Es necesario quitar el
% comentario a alguna de ellas y comentar la instrucción de arriba.
% Se marcan con %----

\documentclass[dat]{fcfmtesis}      

%\makeatletter
\def\vocalc#1{\gdef\@vocalc{#1}}
 \def\@vocalc{\@latex@warning@no@line{No me has dado al \noexpand\vocalb}}
\renewcommand\makeacta{%
             \begin{titlepage}%
             \def\@thanks{}%evita la ocurrencia de doble pie de pagina
  \let \footnotesize \small
  \let \footnoterule \relax
  \let \footnote \thanks %\relax
  %\let \mpfootnotetext{}
  \null\vfill
  %\vskip 60\p@
    \begin{flushleft}%
        {\Large \textbf{T\'{\i}tulo:} \@title \par}
             %\let\@thanks\relax
        {\Large \textbf{Estudiante:}\textsc{\@author} \par}%

    \end{flushleft}%

    \vspace{1cm}
    {\large
      \begin{center}%

        COMIT\'E
        \vspace{2cm}

        \begin{minipage}{5cm}
          \centering
          \vspace{1.2cm}%
          \rule{5cm}{1pt}%

          \@presidente

          Presidente
        \end{minipage}%
        \hspace{1cm}
        \begin{minipage}{5cm}
          \centering
          \vspace{1.2cm}%
          \rule{5cm}{1pt}%

          \@secretario

          Secretario
        \end{minipage}\\*[1cm]

        \begin{minipage}{5cm}
          \centering
          \vspace{1.2cm}%
          \rule{5cm}{1pt}%

          \@vocala

          Vocal
        \end{minipage}%
        \hspace{1cm}
        \begin{minipage}{5cm}
          \centering
          \vspace{1.2cm}%
          \rule{5cm}{1pt}%
        
          \@vocalb
        
          Vocal
        \end{minipage}\\*[1cm]
        
        \begin{minipage}{5cm}
          \centering
          \vspace{1.2cm}%
          \rule{5cm}{1pt}%

          \@vocalc

          Vocal
        \end{minipage}%
        \hspace{1cm}
        \begin{minipage}{5cm}
          \centering
          \vspace{1.2cm}%
          \rule{5cm}{1pt}%

          \@asesor

          Asesor
        \end{minipage}
      \end{center}}
             \@thanks\par
      \vfill\null
  % \@thanks\par          
%
 %  {\@thanks\par}%    
  \end{titlepage}%
      \setcounter{footnote}{0}%
%  \global\let\thanks\relax
%  \global\let\makeacta\relax
%  \global\let\@thanks\@empty
%  \global\let\@author\@empty
%  \global\let\@date\@empty
%  \global\let\@title\@empty
%  \global\let\title\relax
%  \global\let\author\relax
%  \global\let\date\relax
%  \global\let\and\relax
%  \gloabal\let\thanks\relax
%  \global\let\footnote\relax
}
\makeatother
  %----para tener tres vocales, define \vocalc


\usepackage[T1]{fontenc}		   %%% fonts y unicode UTF8
\usepackage[utf8]{inputenc}        
\usepackage{amsmath,amsfonts,amssymb,amsthm}
\setlength{\parindent}{0pt}
% Cualquier paquete necesario debe ser incluido aquí.
% Paquetes ya cargados por default en fcfmtesis:
% babel, latexsym, fancyhdr, graphicx y array

\author{Marcos Bautista Muñoa}                  % Nombre del autor
\title{ Estimación de parámerros epidemiológicos
dependientes del tiempo mediante el Filtro de Kalman por Ensamble en el modelo SIRD}         % Título de la tesis
\asesor{Dra. Beatriz Bonilla Capilla}               % Nombre del asesor
\presidente{Dr. Jorge Velzaquez Castro} % Nombre del presidente del jurado
\secretario{Dr. Erique Varela Carlos}       % Nombre del secretario del jurado
\vocala{Dr. Ivan Fuentecilla Carcamo}         % Nombre del 1er vocal del jurado
                   % Nombre del 2do vocal del jurado
%\vocalc{}                         % Nombre del 3er vocal del jurado
                                   % Usarse sólo si se hizo \makeatletter
\def\vocalc#1{\gdef\@vocalc{#1}}
 \def\@vocalc{\@latex@warning@no@line{No me has dado al \noexpand\vocalb}}
\renewcommand\makeacta{%
             \begin{titlepage}%
             \def\@thanks{}%evita la ocurrencia de doble pie de pagina
  \let \footnotesize \small
  \let \footnoterule \relax
  \let \footnote \thanks %\relax
  %\let \mpfootnotetext{}
  \null\vfill
  %\vskip 60\p@
    \begin{flushleft}%
        {\Large \textbf{T\'{\i}tulo:} \@title \par}
             %\let\@thanks\relax
        {\Large \textbf{Estudiante:}\textsc{\@author} \par}%

    \end{flushleft}%

    \vspace{1cm}
    {\large
      \begin{center}%

        COMIT\'E
        \vspace{2cm}

        \begin{minipage}{5cm}
          \centering
          \vspace{1.2cm}%
          \rule{5cm}{1pt}%

          \@presidente

          Presidente
        \end{minipage}%
        \hspace{1cm}
        \begin{minipage}{5cm}
          \centering
          \vspace{1.2cm}%
          \rule{5cm}{1pt}%

          \@secretario

          Secretario
        \end{minipage}\\*[1cm]

        \begin{minipage}{5cm}
          \centering
          \vspace{1.2cm}%
          \rule{5cm}{1pt}%

          \@vocala

          Vocal
        \end{minipage}%
        \hspace{1cm}
        \begin{minipage}{5cm}
          \centering
          \vspace{1.2cm}%
          \rule{5cm}{1pt}%
        
          \@vocalb
        
          Vocal
        \end{minipage}\\*[1cm]
        
        \begin{minipage}{5cm}
          \centering
          \vspace{1.2cm}%
          \rule{5cm}{1pt}%

          \@vocalc

          Vocal
        \end{minipage}%
        \hspace{1cm}
        \begin{minipage}{5cm}
          \centering
          \vspace{1.2cm}%
          \rule{5cm}{1pt}%

          \@asesor

          Asesor
        \end{minipage}
      \end{center}}
             \@thanks\par
      \vfill\null
  % \@thanks\par          
%
 %  {\@thanks\par}%    
  \end{titlepage}%
      \setcounter{footnote}{0}%
%  \global\let\thanks\relax
%  \global\let\makeacta\relax
%  \global\let\@thanks\@empty
%  \global\let\@author\@empty
%  \global\let\@date\@empty
%  \global\let\@title\@empty
%  \global\let\title\relax
%  \global\let\author\relax
%  \global\let\date\relax
%  \global\let\and\relax
%  \gloabal\let\thanks\relax
%  \global\let\footnote\relax
}
\makeatother


\date{\today}                      % Fecha del examen profesional.
                                   % Escribirse en el formato mes año:
                                   % por ejemplo \date{Julio de 2018}

\hombre

\begin{document}
\frontmatter
           \portada         %----quitar comentario en la versión final
           \maketitle       %----quitar comentario en la versión final
           \makeacta        %----quitar comentario en la versión final
           \setcounter{page}{4}		% La primer página con número impreso
           							% debe ser la siguiente al acta (pág. impar)
           							% pero el conteo se hace a partir del
           							% maketitle, (segunda portada)	
          							% El valor 4 puede cambiarse según corresponda.
           \tableofcontents %----quitar comentario en la versión final

\mainmatter  % A partir de aquí va el trabajo propio, se hará colocando
             % comandos \include{archivo} donde archivo es el nombre de 
             % un archivo de latex con nombre "archivo.tex"
             % Notar que en sistemas Unix las mayúsculas son diferentes
             % a las minúsculas.
\chapter{Introducción}

En diversas situaciones, cuando se intenta llevar la teoría a la práctica, y sobre todo cuando esta se desarrolla en el mundo físico, es necesario considerar una gran cantidad de situaciones distintas.

Esto ocurre al estudiar una epidemia. Aunque existen muchos modelos matemáticos que describen su evolución y comportamiento, es necesario incorporar cada vez más variables para estimar de forma más precisa la situación real a la que se enfrenta la población.

La pandemia de COVID-19 es un claro ejemplo de ello. Si se tuviera una idea fiable de cómo se desarrolla una epidemia, se podrían diseñar políticas públicas realmente útiles.

Esta es una tarea difícil, ya que en el modelo SIRD (susceptibles, infectados, recuperados y fallecidos) normalmente se toman los parámetros de la epidemia como constantes, aun cuando estos pueden cambiar con el tiempo debido a factores como las medidas sanitarias, el comportamiento de la población o la aparición de nuevas variantes.

\section{Objetivos}

Estimar los parámetros $\beta$, $\gamma$ y $\delta$, (tasa de contagio, tasa de recuperación y tasa de mortalidad, respectivamente) dependientes del tiempo, del modelo SIRD mediante el método de ensamble de los filtros de Kalman, a partir de datos sobre COVID-19 en México.

\section{Justificación}

Los filtros de Kalman de ensamble (EnKF, por sus siglas en inglés) son una técnica de asimilación de datos diseñada para estimar estados y parámetros en sistemas dinámicos, especialmente cuando estos son no lineales y están sujetos a incertidumbre. A diferencia del filtro de Kalman clásico, el EnKF utiliza un conjunto de simulaciones (ensamble) para representar la distribución de probabilidad del sistema, lo que permite capturar de manera más flexible la variabilidad y el ruido en los datos. Este enfoque combina las predicciones del modelo con observaciones reales mediante actualizaciones secuenciales, ajustando continuamente las estimaciones conforme se dispone de nueva información.


\chapter{Cadenas de Markov}

Para llegar al procedimiento de Filtros de Ensambles de Kalman, observemos antes la teoría básica y cómo llegamos hasta ese tema. Definamos nuestro espacio de trabajo y desarrollemos paso a paso.

En un tiempo discreto, consideremos un espacio contable $I$. Llamemos a cada $i \in I$ un estado, por lo que $I$ será el espacio de estados. Además, podemos realizar una medición para cada $i \in I$. Entonces, $\lambda = (\lambda_i : i \in I)$ es una medida en $I$ si $0 \leq \lambda_i < \infty$ para todo $i \in I$. Si además la masa total cumple $\sum_{i \in I} \lambda_i = 1$, llamamos a $\lambda$ una distribución. Todo esto lo trabajamos a través de un espacio de probabilidad $(\Omega, \mathcal{F}, \mathbb{P})$.

Si $X$ es una variable aleatoria con valores en $I$, entonces es una función $X: \Omega \rightarrow I$. Suponemos que
\[
\lambda_i = \mathbb{P}(X = i) = \mathbb{P}(\{\omega \in \Omega : X(\omega) = i\}),
\]
lo que nos dice que $\lambda$ define la distribución de $X$. Así, $X$ modela un estado aleatorio que toma el valor $i$ con probabilidad $\lambda_i$.

Ahora, si definimos una matriz $P = (p_{ij} : i, j \in I)$, decimos que es estocástica si cada fila es una distribución. Algunas veces, una representación gráfica de esta matriz ayuda a entender mejor lo que estamos definiendo. Un ejemplo puede verse en la siguiente figura:

% figura

Con lo definido anteriormente, enunciamos la definición de cadena de Markov.

\textbf{Definición.} Decimos que $(X_n)_{n \geq 0}$ es una cadena de Markov con distribución inicial $\lambda$ y matriz de transición $P$ si:
\begin{enumerate}
    \item[(i)] $X_0$ tiene distribución $\lambda$;
    \item[(ii)] para $n \geq 0$, condicionalmente a que $X_n = i$, $X_{n+1}$ tiene distribución $(p_{ij} : j \in I)$ y es independiente de $X_0, \ldots, X_{n-1}$ \cite{Norris}.
\end{enumerate}

La primera característica define el estado inicial y especifica las probabilidades de que el sistema comience en cada uno de los estados posibles en el tiempo cero.

La segunda característica combina dos propiedades fundamentales: la pérdida de memoria y la homogeneidad temporal.

La pérdida de memoria es una propiedad base de las cadenas de Markov. Indica que la distribución del estado futuro depende únicamente del estado actual y es independiente de toda la historia pasada. De forma matemática,
\[
\mathbb{P}(X_{n+1} = j \mid X_n = i, X_{n-1} = i_{n-1}, \ldots, X_0 = i_0) = \mathbb{P}(X_{n+1} = j \mid X_n = i) = p_{ij}.
\]

La homogeneidad temporal se refiere al elemento $p_{ij}$, que indica la probabilidad de transitar del estado $i$ al estado $j$. Esta probabilidad depende únicamente de la pareja de estados $(i, j)$ y no del instante $n$ en el que nos encontremos.
\chapter{Procesos estocásticos y propiedad de Markov (Marco Teorico)}

/section{Procesos estocastcos}

%%%%%%%%%%%%%%%%%%%%%

\section{cadenas de Markov}
Para llegar al procedimiento de Filtros de Ensambles de Kalman, observemos antes la teoría básica y cómo llegamos hasta ese tema. Definamos nuestro espacio de trabajo y desarrollemos paso a paso.

En un tiempo discreto, consideremos un espacio contable $I$. Llamemos a cada $i \in I$ un estado, por lo que $I$ será el espacio de estados. Además, podemos realizar una medición para cada $i \in I$. Entonces, $\lambda = (\lambda_i : i \in I)$ es una medida en $I$ si $0 \leq \lambda_i < \infty$ para todo $i \in I$. Si además la masa total cumple $\sum_{i \in I} \lambda_i = 1$, llamamos a $\lambda$ una distribución. Todo esto lo trabajamos a través de un espacio de probabilidad $(\Omega, \mathcal{F}, \mathbb{P})$.

Si $X$ es una variable aleatoria con valores en $I$, entonces es una función $X: \Omega \rightarrow I$. Suponemos que
\[
\lambda_i = \mathbb{P}(X = i) = \mathbb{P}(\{\omega \in \Omega : X(\omega) = i\}),
\]
lo que nos dice que $\lambda$ define la distribución de $X$. Así, $X$ modela un estado aleatorio que toma el valor $i$ con probabilidad $\lambda_i$.

Ahora, si definimos una matriz $P = (p_{ij} : i, j \in I)$, decimos que es estocástica si cada fila es una distribución. Algunas veces, una representación gráfica de esta matriz ayuda a entender mejor lo que estamos definiendo. Un ejemplo puede verse en la siguiente figura:

% figura

Con lo definido anteriormente, enunciamos la definición de cadena de Markov.

\textbf{Definición.} Decimos que $(X_n)_{n \geq 0}$ es una cadena de Markov con distribución inicial $\lambda$ y matriz de transición $P$ si:
\begin{enumerate}
    \item[(i)] $X_0$ tiene distribución $\lambda$;
    \item[(ii)] para $n \geq 0$, condicionalmente a que $X_n = i$, $X_{n+1}$ tiene distribución $(p_{ij} : j \in I)$ y es independiente de $X_0, \ldots, X_{n-1}$ \cite{Norris}.
\end{enumerate}

La primera característica define el estado inicial y especifica las probabilidades de que el sistema comience en cada uno de los estados posibles en el tiempo cero.

La segunda característica combina dos propiedades fundamentales: la pérdida de memoria y la homogeneidad temporal.

La pérdida de memoria es una propiedad base de las cadenas de Markov. Indica que la distribución del estado futuro depende únicamente del estado actual y es independiente de toda la historia pasada. De forma matemática,
\[
\mathbb{P}(X_{n+1} = j \mid X_n = i, X_{n-1} = i_{n-1}, \ldots, X_0 = i_0) = \mathbb{P}(X_{n+1} = j \mid X_n = i) = p_{ij}.
\]

La homogeneidad temporal se refiere al elemento $p_{ij}$, que indica la probabilidad de transitar del estado $i$ al estado $j$. Esta probabilidad depende únicamente de la pareja de estados $(i, j)$ y no del instante $n$ en el que nos encontremos.

%%%%%%%%%%%%%%%%%%%%%%%%%%%%%%
\section{Suavizadores Bayesianos}

\section{Filtros de kalman}

\section{Filtros de kalman por conjuntos}

\appendix % A partir de esta línea van los apéndices. Quitar si no habrá apéndices
\chapter{SPS Apéndice}
wdadsdsadsdsadsadasdasdsadasdsadasdasdsadasdada
            
Los árticulos ácentúados son mi fuerte!
\begin{figure}[!h]
\centering
\includegraphics[width=.2\textwidth,keepaspectratio=true]{escudo-2}
\end{figure}

mas texto Loren itsu
\[ \vec E = \vec \nabla \cdot \vec \nabla \vec r \times \vec A
\]


\backmatter  % En esta línea está terminado el trabajo. No quitar
\addcontentsline{toc}{chapter}{Bibliografía} % Agregar Bibliografía al índice
\bibliographystyle{plain}
\bibliography{tesis}      % Para utilizar con BibTeX y con un archivo tesis.bib
%\begin{thebibliography}{XX} % Por si alguien no quiere utilizar BibTeX
%\end{thebibliography}
\end{document}            % Fin del documento. cualquier texto a partir de esta
                          % línea es ignorado por latex.
